\documentclass[11pt]{letter}
\usepackage[margin=2.2cm]{geometry}
\usepackage{setspace}
\usepackage{parskip}
\usepackage[hidelinks]{hyperref}
\setstretch{1.05}
\pagestyle{empty}

\signature{January G. Msemakweli \\ On behalf of all co-authors}
\address{January G. Msemakweli \\ Corresponding author \\ Email: jmsemak1@jh.edu}

\begin{document}

\begin{letter}{The Editors \\ \textit{Toxicology} \\ Elsevier}

\opening{Dear Editors,}

We are pleased to submit our manuscript, ``Directional transfer of
chemical-gene effects from rodents to humans: an independence-guarded
meta-analysis of the Comparative Toxicogenomics Database,'' for consideration as
a Research Article in \textit{Toxicology}.

Animal-to-human extrapolation is the central validity problem in toxicology, and
the field is now pivoting toward human-relevant new approach methodologies. Yet
a basic quantity has never been measured at scale: when a chemical is curated to
increase or decrease a given gene in a rodent, how often does the same
directional effect hold in humans? Using the Comparative Toxicogenomics
Database, we assembled 1{,}165{,}406 curated chemical-gene-endpoint effects
across human, mouse and rat and answered this question directly, comparing each
rodent directional call with the directly curated human call for the same
chemical and gene.

We believe the study is a good fit for the journal's emphasis on human-relevant,
mechanistic toxicology for the following reasons.

\begin{itemize}
\item \textbf{A measured human comparator, not an extrapolation.} The primary
estimand compares rodent effects with directly curated human effects, so human
relevance is built into the design rather than assumed.
\item \textbf{An independence guard against circular evidence.} Because a naive
cross-species comparison is inflated when both calls originate from the same
publication, our primary analysis is restricted to effects supported by disjoint
PubMed identifiers across species. This lowers apparent agreement from about
89\% to about 62\%, only marginally above the mouse-versus-rat benchmark, and
reframes cross-species divergence as a problem of reproducibility and context
dependence rather than phylogenetic distance.
\item \textbf{A pre-registered external anchor.} Using orthology and graded
protein sequence conservation from resources independent of the database, we
show that sequence conservation is not associated with directional transfer,
addressing a long-standing question about the ortholog conjecture.
\item \textbf{Actionable framing for evidence integration.} Transfer is governed
by molecular endpoint and replication depth rather than by chemical class,
suggesting that cross-species evidence should be weighted by what was measured
and how firmly it was established.
\end{itemize}

We confirm that this manuscript is original, has not been published previously,
and is not under consideration for publication elsewhere. All authors have
approved the submission and agree to be accountable for the work. The study uses
only openly available, non-identifiable curated data and secondary resources; it
did not involve new experiments in humans or animals. As the analysis is a
meta-analysis of curated literature rather than a controlled exposure study, we
have stated explicitly in the Materials and Methods that individual exposure
doses, durations and chemical identifiers are properties of the underlying source
experiments and are not standardized at the level of this synthesis; our
estimands are therefore literature-level directional concordance conditional on
curation.

The authors declare no competing interests, and the research received no specific
external funding. We have provided highlights, a title page, and comprehensive
supplementary material, and the manuscript is prepared in \LaTeX{} using the
\texttt{elsarticle} class with author-year (Harvard) referencing, as accepted by
the journal.

Thank you for considering our work. We look forward to the reviewers' comments.

\closing{Sincerely,}

\end{letter}
\end{document}
